\section{Move one image through three representations}
A file contains a sequence of bytes. The bitmap display reads words from simulator
memory. A pixel is an interpretation of some of those bytes or words. This project
connects the three: read a binary PPM image, recover its dimensions and pixel samples,
place those pixels in the framebuffer, then write an image that another program can open.

Keep the transformations separate while developing them. First identify where the
header stops. Then check how much raster data must follow. Convert one RGB triplet to
one display word before introducing a loop over the image. Finally, serialize the
words back to RGB bytes and put a new header before them. A wrong header offset and a
wrong component order can both produce a strange image; examining these boundaries
separately tells you which calculation to repair.

Use a two-by-two image with red and green on the first row, blue and white on the
second. Its twelve sample bytes are short enough to inspect individually. The colors
also expose a red/blue reversal and a row-order reversal without relying on a photograph
whose intended appearance may be unclear.

The helper in \path{drafts/ppm-fixtures} writes this initial file for inspection. Run
\texttt{python3 fixture.py} with a new output path. It refuses to overwrite an existing
file. This helper prepares the input; the assembly routines will perform the image
parsing and writing developed in the subsequent construction.
